% sec6_8.tex — Section 6.8: Application: Forward Exchange Rates as Optimal Predictors

In foreign exchange markets, the percentage difference between the forward
exchange rate and the spot exchange rate is called the \textbf{forward premium}. The
relation between the forward premium and the expected rate of currency appreciation/depreciation
over the life of the forward contract has been the subject of a
large number of empirical studies. The simplest hypothesis about the relationship
is that the two are the same. In this section, we test this hypothesis using the same
weekly data set used by Bekaert and Hodrick (1993).

The data cover the period of 1975--1989 and include the following variables:
\begin{align*}
  S_t &= \text{spot exchange rate, stated in units of the foreign currency per dollar} \\
  &\quad \text{(e.g., 125 yen per dollar) on the Friday of week $t$,} \\
  F_t &= \text{30-day forward exchange rate on Friday of week $t$, namely, the price} \\
  &\quad \text{in foreign currency units of the dollar deliverable 30 days from Friday} \\
  &\quad \text{of week $t$,} \\
  S30_t &= \text{the spot rate on the delivery date on a 30-day forward contract made} \\
  &\quad \text{on Friday of week $t$.}
\end{align*}
The important feature of the data is that the maturity of the contract covers several
sampling intervals; the delivery date is after the Friday of week $t + 4$ but before the
Friday of week $t + 5$.

\subsection*{The Market Efficiency Hypothesis}

To indicate the assumptions underlying the hypothesis, consider three strategies
for investing dollar funds for 30~days at date $t$. The first is to invest in a domestic
money market instrument (e.g., U.S.\ Treasury bills) whose maturity is 30~days. Let
$i_t$ be the rate of return over the 30-day period. The second investment strategy is to
convert dollars into foreign currency units, buy a 30-day money market instrument
denominated in the foreign currency whose interest rate is denoted $i_t^*$, and convert
back to dollars at maturity. The rate of return from this investment strategy is
\[
  S_t \cdot (1 + i_t^*) / S30_t.
\]
This strategy involves foreign exchange risk because at time $t$ of investing you
do not know $S30_t$, which is the spot rate 30~days hence. The third investment
strategy differs from the second in that it hedges against exchange risk by buying
dollars forward at price $F_t$ at the same time the dollar fund is converted into foreign
currency units. This hedged foreign investment provides a sure return equal to
\[
  S_t \cdot (1 + i_t^*) / F_t.
\]
Because the first and the third strategies involve no risk, arbitrage requires that the
rates of return be the same:
\[
  1 + i_t = S_t \cdot (1 + i_t^*) / F_t.
\]
Taking logs of both sides and using the approximation $\log(1 + x) \approx x$, we obtain
the \textbf{covered interest parity} equation:
\begin{equation}
  i_t = s_t - f_t + i_t^* \quad \text{or} \quad
  f_t - s_t = i_t^* - i_t,
  \label{eq:6.8.1}
\end{equation}
where lower case letters for $S$ and $F$ are natural logs. That is, the forward premium,
$f_t - s_t$, equals the interest rate differential. This relationship holds almost exactly
in real data. In fact, people routinely calculate the forward premium as the interest
rate differential.

The rate of return from the second investment strategy is uncertain, but if investors
are risk-neutral, its expected return is equal to the sure rate of return. Furthermore,
if their expectations are rational, the expected return is the conditional
expectation based on the information currently available, so
\begin{equation}
  \E[S_t \cdot (1 + i_t^*) / S30_t \mid I_t] = S_t \cdot (1 + i_t^*) / F_t,
  \label{eq:6.8.2}
\end{equation}
where $\E(\cdot \mid I_t)$ is the expectations operator and $I_t$ is the information set as of date
$t$. Since $S_t$, $F_t$, and $i_t^*$ are known at date $t$, this equality reduces to
\[
  \E(F_t / S30_t \mid I_t) = 1.
\]
Again, by log approximation, $F_t / S30_t \approx 1 + f_t - s30_t$. Therefore,
\begin{equation}
  \E(s30_t \mid I_t) = f_t \quad \text{or} \quad
  \E(\varepsilon_t \mid I_t) = 0 \quad \text{where } \varepsilon_t \equiv s30_t - f_t.
  \label{eq:6.8.3}
\end{equation}
This is the market efficiency hypothesis. As is clear from the derivation, the
hypothesis assumes risk-neutrality and rational expectations. The hypothesis is
sometimes described as the forward rate being the optimal forecast of future spot
rates, because the conditional expectation is the optimal forecast in that it minimizes
the mean square error (see Proposition~2.7).

\subsection*{Testing Whether the Unconditional Mean Is Zero}

Taking the unconditional expectation of both sides of \eqref{eq:6.8.3}, we obtain the unconditional
relationship:
\begin{equation}
  \E(s30_t) = \E(f_t) \quad \text{or} \quad \E(\varepsilon_t) = 0.
  \label{eq:6.8.4}
\end{equation}
That is, the forecast must be correct on average. Figure~6.1 is a plot of the forecast
error $\varepsilon_t$ for the yen/dollar exchange rate. It looks like a stationary series, and we
proceed under the assumption that it is.\footnote{Using covered interest rate parity, $\varepsilon_t = (s30_t - s_t) - (i_t^* - i_t)$. If you believe that both the 30-day rate of change $s30_t - s_t$ and the interest rate differential $i_t^* - i_t$ are stationary, you have to believe that the forecast error is stationary.}
Because the forecast error is observable,
we can test the hypothesis that its unconditional mean is zero. This is in contrast
to the Fama exercise, where the forecast error (about future inflation) is observable
only up to a constant.

\begin{figure}[H]
\centering
\includegraphics[width=0.9\textwidth,trim=50 350 50 350,clip]{figures/fig_6_1.png}
\caption{Forecast Error, Yen/Dollar}
\label{fig:6.1}
\end{figure}

To test the hypothesis of market efficiency, we need to derive the asymptotic
distribution of the sample mean, $\bar{\varepsilon}$, under the hypothesis. The task is somewhat
involved because $\{\varepsilon_t\}$ is serially correlated. To see why the forecast error has serial
correlation, note that $I_t$, while including $\{\varepsilon_{t-5}, \varepsilon_{t-6}, \ldots\}$, does not include
$\{\varepsilon_t, \ldots, \varepsilon_{t-4}\}$ because $\varepsilon_t$, which depends on $s30_t$, cannot be observed until after
$t + 4$. Therefore,
\begin{equation}
  \E(\varepsilon_t \mid \varepsilon_{t-5}, \varepsilon_{t-6}, \ldots) = 0.
  \label{eq:6.8.5}
\end{equation}
It follows that $\Cov(\varepsilon_t, \varepsilon_{t-j}) = 0$ for $j \geq 5$ but not necessarily for $j \leq 4$. In
the Fama exercise, the ex-post real interest rate (which is the inflation forecast
error plus a constant) was serially uncorrelated because the maturity for the interest
rate and the sampling interval were the same. Here, the forecast error has serial
correlation because the maturity of the forward contract (30~days) is longer than
the sampling interval (a week). The estimated correlogram is displayed along with
the two-standard error band in Figure~6.2.\footnote{Although theory implies that the population mean is zero, I subtracted the sample mean in the calculation of sample autocorrelations.}
The standard error is calculated under
the assumption that the forecast error is i.i.d. So it equals $1/\sqrt{n}$, as in Table~2.1.
The pattern of autocorrelation is broadly consistent with the hypothesis: it stays
well above two standard errors for the first four lags and generally lies within the
band thereafter.

\begin{figure}[H]
\centering
\includegraphics[width=0.9\textwidth,trim=50 350 50 350,clip]{figures/fig_6_2.png}
\caption{Correlogram of $s30 - f$, Yen/Dollar}
\label{fig:6.2}
\end{figure}

Under the null of market efficiency, since serial correlation vanishes after a
finite lag, it is easy to see that $\{\varepsilon_t\}$ satisfies the essential part of Gordin's condition
restricting serial correlation. By the Law of Iterated Expectations, \eqref{eq:6.8.5} implies
$\E(\varepsilon_t \mid \varepsilon_{t-j}, \varepsilon_{t-j-1}, \ldots) = 0$ for $j \geq 5$, which immediately implies part~(b) of
Gordin's condition. Since the revision of expectations, $r_{tj}$, is zero for $j \geq 5$, part
(c) is satisfied, provided that $r_{tj}$ ($0 \leq j \leq 4$) have finite second moments. It then
follows from Proposition~6.10 that
\begin{equation}
  \frac{\sqrt{n}\,\bar{\varepsilon}}{\sqrt{\gamma_0 + 2 \displaystyle\sum_{j=1}^{4} \gamma_j}}
  = \frac{\bar{\varepsilon}}{\sqrt{\bigl(\gamma_0 + 2 \displaystyle\sum_{j=1}^{4} \gamma_j\bigr)/n}}
  \dto N(0, 1).
  \label{eq:6.8.6}
\end{equation}
Furthermore, the denominator is consistently estimated by replacing $\gamma_j$ by its sample
counterpart $\hat{\gamma}_j$. So the ratio
\begin{equation}
  \frac{\bar{\varepsilon}}{\sqrt{\Bigl(\hat{\gamma}_0 + 2 \displaystyle\sum_{j=1}^{4} \hat{\gamma}_j\Bigr)/n}}
  \label{eq:6.8.7}
\end{equation}
is also asymptotically standard normal. The denominator of this expression will be
called the standard error of the sample mean.

\begin{table}[H]
\centering
\caption{Means of Rates of Change, Forward Premiums,
and Differences between the Two: 1975--1989}
\label{tab:6.1}
\smallskip
\begin{tabular}{lccc}
\toprule
& \multicolumn{3}{c}{Means and Standard Deviations} \\
\cmidrule(lr){2-4}
Exchange & $s30 - s$ & $f - s$ & Difference \\
rate & (actual rate & (expected rate & (unexpected rate \\
& of change) & of change) & of change) \\
\midrule
\textyen/\$ & $-4.98$ & $-3.74$ & $-1.25$ \\
& $(41.6)$ & $(3.65)$ & $(42.4)$ \\
& & & standard error $= 3.56$ \\[4pt]
DM/\$ & $-1.78$ & $-3.91$ & $2.13$ \\
& $(40.6)$ & $(2.17)$ & $(41.1)$ \\
& & & standard error $= 3.26$ \\[4pt]
\pounds/\$ & $3.59$ & $2.16$ & $1.43$ \\
& $(39.2)$ & $(3.49)$ & $(39.9)$ \\
& & & standard error $= 3.29$ \\
\bottomrule
\end{tabular}

\smallskip
\footnotesize
\textsc{Note:} $s30 - s$ is the rate of change over 30~days and $f - s$ is the forward
premium, expressed as annualized percentage changes. Standard deviations are in
parentheses. The standard errors are the values of the denominator of \eqref{eq:6.8.7}. The
data are weekly, from 1975 to 1989. The sample size is 778.
\end{table}

Table~6.1 reports results for testing the implication that the
unconditional mean of the forecast error is zero, for three foreign exchange rates
against the dollar. The three currencies are the Deutsche mark, the British pound,
and the Japanese yen. Columns~1 and~2 of the table report the mean and the standard
deviation of the actual rate of change of the spot rate, $s30_t - s_t$, and the
forward premium, $f_t - s_t$. (Because the spot rate is per dollar, a positive rate of
change means that the dollar strengthened and the
forward premium have been multiplied by 1,200 to express the rates in annualized
percentages. The market efficiency hypothesis is that the forward premium is the
expected rate of change of the spot rate. It is evident that the actual rate of change
is much more volatile than the forward premium. The third column reports the
sample mean of $\varepsilon_t$, which equals the difference between the first two columns for
the sample mean. The numbers in parentheses below the sample mean are the standard
errors calculated as described above. In no case is there significant evidence
against the null hypothesis of zero mean.

\subsection*{Regression Tests}

The unconditional test, however, does not exploit the implication of market efficiency
that the forecast error is orthogonal to conditioning information $I_t$. Probably
the most important variable in the information set $I_t$ for future spot rates is
the forward rate $f_t$. We can test whether the forecast error is uncorrelated with the
forward rate by regressing $\varepsilon_t$ on the constant and $f_t$. This is equivalent to running
the following regression
\begin{equation}
  s30_t = \beta_0 + \beta_1 f_t + \varepsilon_t,
  \label{eq:6.8.8}
\end{equation}
and testing whether $\beta_0 = 0$ and $\beta_1 = 1$. Because the error term is the forecast error
orthogonal to anything known at date $t$ including $f_t$, the regressors are guaranteed
to be orthogonal to the error term. So it might appear that we can estimate \eqref{eq:6.8.8}
by OLS and adjust for serial correlation in the error term as indicated in Section~6.6
for correct inference.

However, as the graph in Figure~6.3 suggests, the spot exchange rate looks like
a process with increasing variance. As we will verify in Chapter~9, we cannot
reject the hypothesis that the process has a unit root. The forward rate has the same
property. (Its plot is not shown here because it is indistinguishable from the plot
of the spot rate.) Thus, Assumption~3.2, which requires the regressor ($f_t$) and the
dependent variable ($s30_t$) to be stationary, is not satisfied. The problem actually
runs deeper. We have observed in Figure~6.1 that the \emph{difference} between $s30_t$ and
$f_t$ looked like a stationary series. In the language to be introduced in Chapter~10,
$s30_t$ and $f_t$ are ``cointegrated'' with a cointegrating vector of $(1, -1)$, in which case
the OLS estimate of $\beta_1$ in \eqref{eq:6.8.8} converges quickly to~1. This can be verified in
Figure~6.4 where $s30_t$ is plotted against $f_t$. If a regression line is fitted, its slope
is 0.99. The problem from the point of view of testing market efficiency is that the
OLS estimate converges to~1 \emph{regardless} of whether or not the hypothesis is true.
Therefore, the test has no power.

\begin{figure}[H]
\centering
\includegraphics[width=0.9\textwidth,trim=50 350 50 350,clip]{figures/fig_6_3.png}
\caption{Yen/Dollar Spot Rate, Jan.\ 1975--Dec.\ 1989}
\label{fig:6.3}
\end{figure}

\begin{figure}[H]
\centering
\includegraphics[width=0.9\textwidth,trim=50 350 50 350,clip]{figures/fig_6_4.png}
\caption{Plot of $s30$ against $f$, Yen/Dollar}
\label{fig:6.4}
\end{figure}

The problem can be dealt with by estimating a different regression,
\begin{equation}
  s30_t - s_t = \beta_0 + \beta_1 (f_t - s_t) + \varepsilon_t.
  \label{eq:6.8.9}
\end{equation}
Judging from the plot of $s30_t - s_t$ and $f_t - s_t$ (not shown), it is reasonable to
assume that they are stationary, which is consistent with Assumption~3.2 (ergodic
stationarity). That the equation satisfies other GMM assumptions under market
efficiency can be verified as follows. For $\beta_0 = 0$ and $\beta_1 = 1$, the error term $\varepsilon_t$
equals the forecast error $s30_t - f_t$ and so the orthogonality conditions $\E(\varepsilon_t) = 0$ and
$\E[(f_t - s_t) \cdot \varepsilon_t] = 0$ (Assumption~3.3 with $\mathbf{z}_t = \mathbf{x}_t$) are satisfied. The identification
condition is that the second moment of $(1, f_t - s_t)$ be nonsingular, which is true if
and only if $\Var(f_t - s_t) > 0$. This is evidently satisfied; if the population variance
were zero, we would not observe any variations in $f_t - s_t$. This leaves us with
Assumption~3.5$'$, which requires that there is not too much serial correlation in
\[
  \mathbf{g}_t \equiv \mathbf{x}_t \cdot \varepsilon_t =
  \begin{bmatrix} \varepsilon_t \\ (f_t - s_t)\,\varepsilon_t \end{bmatrix}.
\]
As already noted, since $I_t$ includes $\{\varepsilon_{t-5}, \varepsilon_{t-6}, \ldots\}$, the forecast error satisfies
\eqref{eq:6.8.5}. Review Question~3 will ask you to show that $\{\mathbf{g}_t\}$ inherits the same property,
namely,
\begin{equation}
  \E(\mathbf{g}_t \mid \mathbf{g}_{t-5}, \mathbf{g}_{t-6}, \ldots) = \mathbf{0}.
  \label{eq:6.8.10}
\end{equation}
So, as in the unconditional test, Gordin's condition is satisfied, provided that the
relevant second moments are finite. Since $\E(\mathbf{g}_t \mathbf{g}_{t-j}') = \mathbf{0}$ for $j \geq 5$, the long-run
covariance matrix $\mathbf{S}$ can be estimated as in \eqref{eq:6.6.4} with $q = 4$. (In the empirical
exercise, you will be asked to estimate $\hat{\mathbf{S}}$ by the Bartlett kernel and also by
VARHAC.)

To summarize, the efficient GMM estimator of $\bm{\beta} = (\beta_0, \beta_1)'$ in \eqref{eq:6.8.9} is OLS
and the consistent estimator of its asymptotic variance is given by
\begin{equation}
  \widehat{\avar(\hat{\bm{\beta}}_{\text{OLS}})}
  = \mathbf{S}_{xx}^{-1}\,\hat{\mathbf{S}}\,\mathbf{S}_{xx}^{-1},
  \label{eq:6.8.11}
\end{equation}
where $\hat{\mathbf{S}}$ is given in \eqref{eq:6.6.4} with $q = 4$. Table~6.2 reports the regression test results
for the three currencies. The corresponding plot for the yen/\$ exchange rate is
in Figure~6.5. Notice that the estimated slope coefficient is not only significantly
different from~1 but also negative, for all three currencies. The null hypothesis
implied by market efficiency that $\beta_0 = 0$ and $\beta_1 = 1$ can be rejected decisively.
This is reflected in the too many observations in the north-west quadrant of Figure~6.5;
too often expected dollar depreciation (negative $f_t - s_t$) is associated with
actual dollar \emph{appreciation} (positive $s30_t - s_t$).

\begin{table}[H]
\centering
\caption{Regression Tests of Market Efficiency: 1975--1989}
\label{tab:6.2}
\smallskip
$s30_t - s_t = \beta_0 + \beta_1 (f_t - s_t) + \varepsilon_t$

\medskip
\begin{tabular}{lcccc}
\toprule
& \multicolumn{2}{c}{Regression Coefficients} & & Wald Statistic for \\
\cmidrule(lr){2-3}
Currency & Constant & Forward & $R^2$ & H${}_0$: $\beta_0 = 0$, $\beta_1 = 1$ \\
& & premium & & \\
\midrule
\textyen/\$ & $-12.8$ & $-2.10$ & 0.034 & 18.6 \\
& $(4.01)$ & $(0.738)$ & & $(p = 0.009\%)$ \\[4pt]
DM/\$ & $-13.6$ & $-3.01$ & 0.026 & 8.7 \\
& $(5.72)$ & $(1.37)$ & & $(p = 1.312\%)$ \\[4pt]
\pounds/\$ & $7.96$ & $-2.02$ & 0.033 & 12.9 \\
& $(3.54)$ & $(0.85)$ & & $(p = 0.156\%)$ \\
\bottomrule
\end{tabular}

\smallskip
\footnotesize
\textsc{Note:} Standard errors are in parentheses. They are heteroskedasticity-robust and
allow for serial correlation up to four lags calculated by the formula \eqref{eq:6.8.11}. The
sample size is~778. Our results differ slightly from those in Bekaert and Hodrick
(1993) because they used the Bartlett kernel-based estimator with $q(n) = 4$ to
estimate $\hat{\mathbf{S}}$.
\end{table}

\begin{figure}[H]
\centering
\includegraphics[width=0.9\textwidth,trim=50 350 50 350,clip]{figures/fig_6_5.png}
\caption{Plot of $s30 - s$ against $f - s$, Yen/Dollar}
\label{fig:6.5}
\end{figure}

\subsection*{Questions for Review}

\begin{enumerate}
\item (The standard error in the unconditional test) \; Let $\varepsilon_t = s30_t - f_t$ and consider
a regression of $\varepsilon_t$ on the constant. To account for serial correlation, use \eqref{eq:6.6.4}
with $q = 4$ for $\hat{\mathbf{S}}$. Verify that the $t$-value for the hypothesis that the intercept is
zero is given by \eqref{eq:6.8.7}.

\item (Lack of identification) \; For simplicity suppose $F_t$ is the two-week forward
exchange rate so that the $s30_t$ in \eqref{eq:6.8.9} can be replaced by $s_{t+2}$. Suppose that
the spot exchange rate is a random walk and $I_t = \{s_t, s_{t-1}, \ldots\}$. What should
$f_t$ be under the hypothesis? Which of the required assumptions are violated?
[Answer: The identification condition.] Suppose, instead, that $s_t$ is AR(1)
satisfying $s_t = c + \phi s_{t-1} + \eta_t$ where $\{\eta_t\}$ is i.i.d.\ and $|\phi| < 1$. Verify that
$f_t = (1 + \phi) c + \phi^2 s_t$. \textbf{Hint:} $s_{t+2} = (1 + \phi) c + \phi^2 s_t + (\eta_{t+2} + \phi \eta_{t+1})$. Verify
that $\varepsilon_t = \eta_{t+2} + \phi \eta_{t+1}$.

\item Verify \eqref{eq:6.8.10}. \textbf{Hint:} $\E(\mathbf{g}_t \mid \mathbf{g}_{t-5}, \mathbf{g}_{t-6}, \ldots) = \E[\E(\mathbf{g}_t \mid I_t) \mid \mathbf{g}_{t-5}, \ldots]$
by the Law of Iterated Expectations, since $(f_t - s_t)$ and $\varepsilon_t$ are both in $I_t$.
\end{enumerate}
